\section*{第 \ref{ch:g12:equilibrium} 章——化学平衡：商与常数}

\begin{solution}{exo:g12:equilibrium:1}
(a) $Q = \dfrac{[\ce{NH3}]^2 (c^\circ)^2}{[\ce{N2}][\ce{H2}]^3}$;
(b) $Q = \dfrac{[\ce{CH3COO-}][\ce{H3O+}]}{[\ce{CH3COOH}]\,c^\circ}$ （水为
\omterm{def:g6:solutions-and-solubility:solute}{溶剂}，不计入）；
(c) $Q = [\ce{Ca^{2+}}][\ce{CO3^{2-}}]/(c^\circ)^2$ （固体不计入）；
(d) $Q = \dfrac{[\ce{FeSCN^{2+}}]\,c^\circ}{[\ce{Fe^{3+}}][\ce{SCN-}]}$.
\end{solution}

\begin{solution}{exo:g12:equilibrium:2}
$\tau = 0.020/0.050 = 0.40$：\omterm{def:g12:equilibrium:non-total}{非完全反应}。
\end{solution}

\begin{solution}{exo:g12:equilibrium:3}
$Q = 2 < K$：正向。$Q = 50 > K$：逆向。$Q = K$：平衡不移动。
\end{solution}

\begin{solution}{exo:g12:equilibrium:4}
组成保持不变，但正、逆反应仍在进行，
且速率相等：\omterm{def:g7:atoms-and-molecules:molecule}{分子}仍在两个方向上不断反应。
\end{solution}

\begin{solution}{exo:g12:equilibrium:5}
否，$K$ 不变（\omterm{def:g10:reaction-progress-table:system}{最终状态}相同）；达到
平衡所需的时间更短。
\end{solution}

\begin{solution}{exo:g12:equilibrium:6}
$\dfrac{x^2}{(2-x)^2} = 4$，故 $\dfrac{x}{2-x} = 2$，且
$x = \qty{1.33}{mol}$：酸与醇各 \qty{0.67}{mol}，
酯与水各 \qty{1.33}{mol}（同样为三分之二）。
\end{solution}

\begin{solution}{exo:g12:equilibrium:7}
$Q = (0.5 \times 0.5)/(0.5 \times 0.5) = 1 < 4$：正向，
生成更多的酯。
\end{solution}

\begin{solution}{exo:g12:equilibrium:8}
由酸一侧读得约 \qty{0.52}{mol}，由
酯一侧读得约 \qty{0.79}{mol}。约 \qty{100}{h} 后，两条曲线均落在虚线上，
即达到平衡。
\end{solution}

\begin{solution}{exo:g12:equilibrium:9}
设已水解 $y$ \omterm{def:g10:the-mole:amount}{摩尔}：$\dfrac{y^2}{(1-y)^2} = \dfrac{1}{4}$，故
$\dfrac{y}{1-y} = \dfrac{1}{2}$，$y = \frac{1}{3}$：酯与水各 $\frac{2}{3}$ \omterm{def:g10:the-mole:amount}{摩尔}，
酸与醇各 $\frac{1}{3}$ \omterm{def:g10:the-mole:amount}{摩尔}。与酯化反应的
最终\omterm{def:g6:pure-substances-and-mixtures:mixture}{混合物}相同。
\end{solution}

\begin{solution}{exo:g12:equilibrium:10}
固体是纯相，其“浓度”不变，因此
不计入 $Q$。加入更多固体不改变 $Q$，所以
也不使平衡移动（只要仍有固体存在）。
\end{solution}

\begin{solution}{exo:g12:equilibrium:11}
$K$ 决定了气体与水中二氧化碳的比值。
打开后，水面上方的气体逸入\omterm{def:g4:air-a-mixture-of-gases:air}{空气}，
气体中的二氧化碳减少，$Q < K$，于是
更多的二氧化碳从水中逸出，直至几乎不剩。
\end{solution}

\begin{solution}{exo:g12:equilibrium:12}
平衡时：酸与醇各 $\frac{1}{3}$，酯与水
各 $\frac{2}{3}$。移去一半的水后，剩余 $\frac{1}{3}$：
$Q = \dfrac{(2/3)(1/3)}{(1/3)(1/3)} = 2 < 4$，正向。若再有 $y$ \omterm{def:g10:the-mole:amount}{摩尔}
参与反应，则 $(2/3 + y)(1/3 + y) = 4(1/3 - y)^2$，即
$27y^2 - 33y + 2 = 0$，$y = 0.064$：酯升至 \qty{0.73}{mol}。
\end{solution}

\begin{solution}{exo:g12:equilibrium:13}
逆反应的\omterm{def:g12:equilibrium:quotient}{反应商}，其\omterm{def:g7:atoms-and-molecules:molecule}{分子}与分母
互换：$Q' = 1/Q$，故平衡时 $K' = 1/K$。水解的
\omterm{def:g12:equilibrium:constant}{平衡常数}为 $1/4 = 0.25$。
\end{solution}

\begin{solution}{exo:g12:equilibrium:14}
$Q = (2 \times 2)/(1 \times 1) = 4 = K$：\omterm{def:g6:pure-substances-and-mixtures:mixture}{混合物}已
处于\omterm{def:g12:equilibrium:equilibrium-state}{平衡状态}，其组成不再改变。
\end{solution}

\begin{solution}{exo:g12:equilibrium:15}
$n - 0.95 = 0.9025 / (4 \times 0.05) = 4.51$，故每\omterm{def:g10:the-mole:amount}{摩尔}酸需
醇 $n = \qty{5.5}{mol}$。
\end{solution}

\begin{solution}{pb:g12:equilibrium:1}
\textbf{1.} \ce{CH3-COOH + CH3-CH2-OH <=> CH3-COO-CH2-CH3 + H2O}.

\textbf{2.} 该反应为\omterm{def:g12:equilibrium:non-total}{非完全反应}：它停在
正、逆两个方向仍在进行的\omterm{def:g12:equilibrium:equilibrium-state}{平衡状态}。

\textbf{3.} $Q = \dfrac{n_{\text{酯}}\, n_{\text{水}}}{n_{\text{酸}}\, n_{\text{醇}}}$。

\textbf{4.} $x_{\max} = \qty{1}{mol}$；$x_f = \frac{2}{3}$ \omterm{def:g10:the-mole:amount}{摩尔}。

\textbf{5.} $\tau = 0.67$.

\textbf{6.} 酸与醇各 $\frac{1}{3}$ \omterm{def:g10:the-mole:amount}{摩尔}，酯与水
各 $\frac{2}{3}$ \omterm{def:g10:the-mole:amount}{摩尔}。

\textbf{7.} $K = \dfrac{(2/3)^2}{(1/3)^2} = 4$.

\textbf{8.} 三分之一发生水解：酯与水各 $\frac{2}{3}$，酸
与醇各 $\frac{1}{3}$：同一\omterm{def:g6:pure-substances-and-mixtures:mixture}{混合物}，$Q_{eq} = 4$。

\textbf{9.} 同一温度下的同一反应，且 $K$
仅取决于温度。

\textbf{10.} 酸 $1 - x$；醇 $3 - x$；酯 $x$；水 $x$。

\textbf{11.} $x^2 = 4(1 - x)(3 - x) = 4(3 - 4x + x^2)$，故
$3x^2 - 16x + 12 = 0$.

\textbf{12.} $x = (16 \pm \sqrt{256 - 144})/6 = (16 \pm 10.58)/6$：0.903
或 4.43；只有 $x_f = \qty{0.903}{mol}$ 介于 0 与 1 之间。

\textbf{13.} \qty{90}{\percent}.

\textbf{14.} $M = \qty{88.0}{g/mol}$: $0.903 \times 88.0 = \qty{79.5}{g}$.

\textbf{15.} 每移去一次水，都会降低 $Q$ 的\omterm{def:g7:atoms-and-molecules:molecule}{分子}，使其
降至 $K$ 以下：体系再次正向移动。

\textbf{16.} 反应将持续进行，直至酸耗尽：
\qty{100}{\percent}.

\textbf{17.} 乙醇更便宜，且其过量部分
易与酯分离（并可\omterm{def:g5:raw-materials-and-recycling:recycling}{回收利用}）。

\textbf{18.} 只影响时间：放热极少的反应，其
\omterm{def:g12:equilibrium:constant}{平衡常数}几乎不随温度变化，升温只会使
平衡更快到达。

\textbf{19.} 约 \qty{90}{\percent} 的酸。
\end{solution}
